
\documentclass{article}
\usepackage{colortbl}
\usepackage{makecell}
\usepackage{multirow}
\usepackage{supertabular}

\begin{document}

\newcounter{utterance}

\centering \large Interaction Transcript for game `taboo', experiment `high\_en', episode 12 with reka{-}flash{-}3.1{-}t0.0.
\vspace{24pt}

{ \footnotesize  \setcounter{utterance}{1}
\setlength{\tabcolsep}{0pt}
\begin{supertabular}{c@{$\;$}|p{.15\linewidth}@{}p{.15\linewidth}p{.15\linewidth}p{.15\linewidth}p{.15\linewidth}p{.15\linewidth}}
    \# & $\;$A & \multicolumn{4}{c}{Game Master} & $\;\:$B\\
    \hline

    \theutterance \stepcounter{utterance}  
    & & \multicolumn{4}{p{0.6\linewidth}}{
        \cellcolor[rgb]{0.9,0.9,0.9}{
            \makecell[{{p{\linewidth}}}]{
                \texttt{\tiny{[P1$\langle$GM]}}
                \texttt{You are playing a collaborative word guessing game in which you have to describe a target word for another player to guess.} \\
\\ 
\texttt{Rules:} \\
\texttt{(a) You have to reply in the form: CLUE: <some text>. Guesses from the other player will start with GUESS.} \\
\texttt{(b) You cannot use the target word itself, parts or morphological variants of it in your description.} \\
\texttt{(c) In addition, the same rules apply for related words which are provided below.} \\
\\ 
\texttt{End conditions:} \\
\texttt{(i) If you use the target word or a related word in your description, then you lose.} \\
\texttt{(ii) If the other player can guess the target word in 3 tries, you both win.} \\
\\ 
\texttt{Let us start.} \\
\\ 
\texttt{This is the target word that you need to describe and that the other player needs to guess:} \\
\\ 
\texttt{wing} \\
\\ 
\texttt{Related words are:} \\
\\ 
\texttt{{-} feather} \\
\texttt{{-} fly} \\
\texttt{{-} bird} \\
\\ 
\texttt{Important: You are under time pressure, give short descriptions that are to the point!} \\
            }
        }
    }
    & \\ \\

    \theutterance \stepcounter{utterance}  
    & \multicolumn{4}{p{0.6\linewidth}}{
        \cellcolor[rgb]{0.9,0.9,0.9}{
            \makecell[{{p{\linewidth}}}]{
                \texttt{\tiny{[P1$\rangle$GM]}}
                \texttt{<reasoning>} \\
\texttt{I need to come up with a clue for the word "wing" without using "wing" itself, "feather," "fly," or "bird." Let's start by thinking about what a wing is. Wings are part of birds, right? But I can't mention bird. So maybe focus on other animals? Like bats have wings, but bat is a related word? The related words given are feather, fly, bird. So bat isn't listed, so maybe it's allowed. But maybe the user wants to avoid that. Alternatively, think about the function of a wing. It's used for flying. But "fly" is a related word, so I can't use that either. So maybe describe something else about wings.} \\
\\ 
\texttt{Wings allow animals to move through the air. But "fly" is off{-}limits. Maybe think about the structure. A wing has bones, muscles, maybe feathers. But "feather" is a related word. So can't use that. What else? Maybe the shape? Like a wing is a flap or a part that's curved. Or maybe something about how it's used. For example, "used to glide" but glide isn't fly. Wait, glide is different from fly. Maybe "glide" is okay. Let me check the rules again. The related words are feather, fly, bird. So glide isn't related, so maybe that's allowed. So maybe "Used to glide through the air." But does that use "fly"? No, glide is different. Alternatively, "Part that lets you soar." Soar is related to flying but not the same as fly. Maybe that's okay. Let me think of other possibilities. What about the opposite of a leg? Like, instead of legs, wings are for movement. But how to phrase that without using the word. Maybe "Opposite of legs for movement in the sky." But "sky" is okay. Wait, but the clue has to be concise. Maybe "Opposite of legs, enables flight." But "flight" is related to fly. Hmm. Alternatively, "Two flaps for movement in the air." But "flaps" might be okay. Wait, the user said "wing" is the target, so parts like flap are allowed as long as they don't use the target word. Let me check the rules again. Rule (b) says you can't use the target word or parts} \\
            }
        }
    }
    & & \\ \\

    \theutterance \stepcounter{utterance}  
    & & & \multicolumn{2}{p{0.3\linewidth}}{
        \cellcolor[rgb]{0.9,0.9,0.9}{
            \makecell[{{p{\linewidth}}}]{
                \texttt{\tiny{[GM$|$GM]}}
                \texttt{abort game} \\
            }
        }
    }
    & & \\ \\

\end{supertabular}
}

\end{document}
